\section{Por tema}
As listas saem do \texttt{.bib}: cada obra traz autores, ano e o link
(cópia aberta, ou o DOI marcado \enquote{acesso pago}).

\subsection{TDD}
Estudos empíricos sobre o que o TDD muda na qualidade, na produtividade e
no design.
\ondeachar*{tema-tdd}

\subsection{Efetividade dos testes: cobertura, mutação e propriedades}
\ondeachar*{tema-efetividade}

\subsection{Testes instáveis, test smells, carga e caos}
\ondeachar*{tema-flaky}

\subsection{Mocks na JVM}
\ondeachar*{tema-mocks}

\subsection{Arquitetura de software}
Não achamos paper revisado por pares sobre arquitetura hexagonal ou Clean
Architecture, nem estudo empírico da pirâmide de testes; estes são a base
conceitual e empírica mais próxima.
\ondeachar*{tema-arquitetura}

\subsection{JIT e otimização adaptativa}
De Self ao HotSpot e ao Graal.
\ondeachar*{tema-jit}

\subsection{Memória: coleta de lixo e modelo de memória}
\ondeachar*{tema-memoria}

\subsection{Metodologia de benchmark}
\ondeachar*{tema-bench}

\subsection{GraalVM, Truffle e Native Image}
\ondeachar*{tema-graal}

\subsection{Outras VMs e modelos históricos}
BEAM, P-code, Smalltalk, Self, CLR, Dalvik e Lua. Não achamos estudo
comparativo de Dalvik e ART em venue de primeira linha, nem paper revisado
por pares sobre threads virtuais (para estas, veja a JEP~444).
\ondeachar*{tema-vms}

\subsection{Cursos universitários abertos}
\ondeachar*{tema-cursos}
